% =============================================================================
% EBF DOCUMENT SPECIFICATION FRAMEWORK
% =============================================================================
%
% Meta-Framework zur vollständigen Dokumenten-Spezifikation
% 
% Aus den 12 Spezifikations-Dimensionen wird automatisch abgeleitet:
% - Dokumenten-Format (1-Pager, Memo, Paper, Buch, etc.)
% - Scope & Struktur
% - Qualitätsanforderungen
%
% =============================================================================

\documentclass[12pt,a4paper]{article}
\usepackage[margin=2.5cm]{geometry}
\usepackage{amsmath,amssymb}
\usepackage{booktabs}
\usepackage{longtable}
\usepackage{tcolorbox}
\usepackage{enumitem}
\usepackage{hyperref}
\usepackage{xcolor}
\usepackage{array}

\title{EBF Document Specification Framework\\
\large Von der Spezifikation zum Format}
\author{BEATRIX Research Group}
\date{Januar 2026}

\begin{document}
\maketitle
\tableofcontents
\newpage

% =============================================================================
\section{Die 12 Spezifikations-Dimensionen}
% =============================================================================

\begin{tcolorbox}[colback=blue!5!white, colframe=blue!75!black, 
    title=Dokumenten-DNA: Vollständige Spezifikation]

Jedes Dokument in EBF wird durch \textbf{12 Dimensionen} vollständig spezifiziert.
Aus diesen Dimensionen lässt sich das optimale Format automatisch ableiten.

\end{tcolorbox}

% -----------------------------------------------------------------------------
\subsection{Dimension 1: ZIEL (Purpose)}
% -----------------------------------------------------------------------------

\textbf{Frage:} Was soll mit dem Dokument erreicht werden?

\begin{table}[htbp]
\centering
\renewcommand{\arraystretch}{1.3}
\begin{tabular}{@{}lp{8cm}l@{}}
\toprule
\textbf{Ziel-Typ} & \textbf{Beschreibung} & \textbf{Format-Tendenz} \\
\midrule
Informieren & Wissen vermitteln, Fakten darstellen & Report, Paper \\
Überzeugen & Meinung ändern, Entscheidung herbeiführen & Memo, Proposal \\
Dokumentieren & Festhalten für spätere Referenz & Appendix, Manual \\
Anleiten & Handlungen ermöglichen & Guide, Tutorial \\
Entscheiden & Entscheidungsgrundlage liefern & Decision Paper \\
Nachweisen & Compliance, Audit, Beleg & Protocol, Log \\
Theoretisieren & Neue Konzepte entwickeln & Paper, Monograph \\
Integrieren & Bestehendes zusammenführen & Review, Synthesis \\
\bottomrule
\end{tabular}
\end{table}

% -----------------------------------------------------------------------------
\subsection{Dimension 2: LESER (Audience)}
% -----------------------------------------------------------------------------

\textbf{Frage:} Wer ist die primäre Zielgruppe?

\begin{table}[htbp]
\centering
\renewcommand{\arraystretch}{1.3}
\begin{tabular}{@{}lp{5cm}p{5cm}@{}}
\toprule
\textbf{Leser-Typ} & \textbf{Charakteristik} & \textbf{Implikation} \\
\midrule
Experten & Tiefes Fachwissen, wenig Zeit & Dicht, formal, keine Basics \\
Fachleute & Grundwissen, mittlere Zeit & Balanciert, erklärt Neues \\
Generalisten & Breites Wissen, variable Zeit & Zugänglich, kontextualisiert \\
Entscheider & Wenig Zeit, braucht Essenz & Kurz, Empfehlungen vorne \\
Studenten & Lernend, viel Zeit & Pädagogisch, Beispiele \\
Öffentlichkeit & Kein Fachwissen & Einfach, visualisiert \\
Regulatoren & Compliance-fokussiert & Präzise, vollständig \\
Peer Reviewer & Kritisch, sucht Fehler & Rigoros, belegt \\
\bottomrule
\end{tabular}
\end{table}

% -----------------------------------------------------------------------------
\subsection{Dimension 3: ANLASS (Trigger)}
% -----------------------------------------------------------------------------

\textbf{Frage:} Warum wird das Dokument jetzt erstellt?

\begin{table}[htbp]
\centering
\renewcommand{\arraystretch}{1.3}
\begin{tabular}{@{}lp{6cm}l@{}}
\toprule
\textbf{Anlass} & \textbf{Beispiel} & \textbf{Dringlichkeit} \\
\midrule
Anfrage & Kunde/Stakeholder fragt & Hoch \\
Deadline & Konferenz, Journal, Projekt & Hoch \\
Meilenstein & Projektphase abgeschlossen & Mittel \\
Erkenntnisgewinn & Neue Einsicht dokumentieren & Mittel \\
Systematisierung & Ordnung schaffen & Niedrig \\
Prävention & Für zukünftigen Bedarf & Niedrig \\
Pflicht & Regulatorisch, vertraglich & Extern definiert \\
Opportunität & Gelegenheit nutzen & Variabel \\
\bottomrule
\end{tabular}
\end{table}

% -----------------------------------------------------------------------------
\subsection{Dimension 4: SCOPE (In-Scope)}
% -----------------------------------------------------------------------------

\textbf{Frage:} Was genau wird behandelt?

\begin{tcolorbox}[colback=green!5!white, colframe=green!75!black]
\textbf{Scope-Definition erfordert:}
\begin{itemize}[nosep]
\item Thematische Grenzen (welche Konzepte)
\item Zeitliche Grenzen (welcher Zeitraum)
\item Räumliche Grenzen (welche Regionen/Märkte)
\item Theoretische Grenzen (welche Annahmen)
\item Empirische Grenzen (welche Daten)
\item Analytische Tiefe (wie detailliert)
\end{itemize}
\end{tcolorbox}

% -----------------------------------------------------------------------------
\subsection{Dimension 5: OUT-OF-SCOPE}
% -----------------------------------------------------------------------------

\textbf{Frage:} Was wird explizit NICHT behandelt?

\begin{tcolorbox}[colback=red!5!white, colframe=red!75!black]
\textbf{Wichtig:} Out-of-Scope muss \textbf{explizit} dokumentiert werden!

\textbf{Typische Out-of-Scope Kategorien:}
\begin{itemize}[nosep]
\item Verwandte aber nicht zentrale Themen
\item Alternative Theorien/Ansätze
\item Implementierungsdetails (wenn konzeptuell)
\item Konzeptuelle Grundlagen (wenn angewandt)
\item Historische Entwicklung (wenn aktuell)
\item Zukünftige Entwicklungen (wenn deskriptiv)
\item Normative Bewertungen (wenn deskriptiv)
\end{itemize}
\end{tcolorbox}

% -----------------------------------------------------------------------------
\subsection{Dimension 6: MINDESTANFORDERUNGEN (Must-Haves)}
% -----------------------------------------------------------------------------

\textbf{Frage:} Was MUSS das Dokument enthalten?

\begin{table}[htbp]
\centering
\renewcommand{\arraystretch}{1.3}
\begin{tabular}{@{}lp{9cm}@{}}
\toprule
\textbf{Kategorie} & \textbf{Beispiel-Anforderungen} \\
\midrule
Inhaltlich & Bestimmte Themen, Definitionen, Ergebnisse \\
Strukturell & Executive Summary, Appendix, Glossar \\
Formal & Zitierweise, Formatierung, Länge \\
Evidenz & Mindestanzahl Quellen, Datentypen \\
Qualitativ & Peer Review, Lektorat \\
Rechtlich & Disclaimer, Copyright, Datenschutz \\
Prozessual & Freigaben, Versionierung \\
\bottomrule
\end{tabular}
\end{table}

% -----------------------------------------------------------------------------
\subsection{Dimension 7: QUALITÄTSSTANDARD}
% -----------------------------------------------------------------------------

\textbf{Frage:} Welches Qualitätsniveau wird erwartet?

\begin{table}[htbp]
\centering
\renewcommand{\arraystretch}{1.3}
\begin{tabular}{@{}lp{4cm}p{4cm}l@{}}
\toprule
\textbf{Level} & \textbf{Charakteristik} & \textbf{Prüfung} & \textbf{Typisch für} \\
\midrule
Q1: Draft & Erste Version, Fehler ok & Selbst-Review & Interne Notizen \\
Q2: Working & Lesbar, grob geprüft & Team-Review & Working Papers \\
Q3: Professional & Poliert, konsistent & Expert Review & Reports \\
Q4: Publication & Peer-reviewed Qualität & External Review & Journals \\
Q5: Reference & Definitive Version & Multi-stage Review & Textbooks \\
\bottomrule
\end{tabular}
\end{table}

% -----------------------------------------------------------------------------
\subsection{Dimension 8: EVIDENZ-ANFORDERUNG}
% -----------------------------------------------------------------------------

\textbf{Frage:} Welche Art von Belegen werden benötigt?

\begin{table}[htbp]
\centering
\renewcommand{\arraystretch}{1.3}
\begin{tabular}{@{}lp{5cm}l@{}}
\toprule
\textbf{Evidenz-Typ} & \textbf{Beschreibung} & \textbf{Stärke} \\
\midrule
Anekdotisch & Einzelbeispiele, Illustrationen & Niedrig \\
Expert Opinion & Aussagen von Autoritäten & Mittel-Niedrig \\
Theoretisch & Logische Ableitungen & Mittel \\
Empirisch-deskriptiv & Beobachtungsdaten & Mittel \\
Empirisch-kausal & RCTs, Natural Experiments & Hoch \\
Meta-analytisch & Systematische Reviews & Sehr hoch \\
Formal-mathematisch & Beweise, Ableitungen & Kontextabhängig \\
\bottomrule
\end{tabular}
\end{table}

% -----------------------------------------------------------------------------
\subsection{Dimension 9: LEBENSZYKLUS}
% -----------------------------------------------------------------------------

\textbf{Frage:} Wie entwickelt sich das Dokument über Zeit?

\begin{table}[htbp]
\centering
\renewcommand{\arraystretch}{1.3}
\begin{tabular}{@{}lp{5cm}p{4cm}@{}}
\toprule
\textbf{Typ} & \textbf{Beschreibung} & \textbf{Implikation} \\
\midrule
Einmalig & Fixiert nach Fertigstellung & Klare Deadline, kein Update \\
Versioniert & Periodische Major Updates & Versionsnummern, Changelog \\
Living Document & Kontinuierlich aktualisiert & Update-Prozess, Ownership \\
Archiviert & Nach Periode eingefroren & Archivierungsdatum \\
Ephemer & Kurze Lebensdauer & Verfallsdatum \\
\bottomrule
\end{tabular}
\end{table}

% -----------------------------------------------------------------------------
\subsection{Dimension 10: VERTRAULICHKEIT}
% -----------------------------------------------------------------------------

\textbf{Frage:} Wer darf das Dokument sehen?

\begin{table}[htbp]
\centering
\renewcommand{\arraystretch}{1.3}
\begin{tabular}{@{}lp{5cm}l@{}}
\toprule
\textbf{Level} & \textbf{Zugang} & \textbf{Beispiel} \\
\midrule
Öffentlich & Jeder & Website, Open Access Paper \\
Kunden & Vertragspartner & Client Reports \\
Intern & Nur Mitarbeiter & Working Papers \\
Restricted & Need-to-Know & Strategiedokumente \\
Confidential & Einzelne Personen & M\&A, Personal \\
\bottomrule
\end{tabular}
\end{table}

% -----------------------------------------------------------------------------
\subsection{Dimension 11: ABHÄNGIGKEITEN}
% -----------------------------------------------------------------------------

\textbf{Frage:} Was muss existieren/passieren, bevor das Dokument erstellt werden kann?

\begin{tcolorbox}[colback=yellow!5!white, colframe=yellow!75!black]
\textbf{Abhängigkeits-Typen:}
\begin{itemize}[nosep]
\item \textbf{Inhaltlich:} Andere Dokumente, Daten, Analysen
\item \textbf{Prozessual:} Freigaben, Reviews, Meetings
\item \textbf{Ressourcen:} Budget, Personal, Tools
\item \textbf{Zeitlich:} Events, Deadlines anderer Projekte
\item \textbf{Extern:} Third-party Inputs, Regulierung
\end{itemize}
\end{tcolorbox}

% -----------------------------------------------------------------------------
\subsection{Dimension 12: LIEFEROBJEKTE}
% -----------------------------------------------------------------------------

\textbf{Frage:} Was genau wird am Ende abgeliefert?

\begin{table}[htbp]
\centering
\renewcommand{\arraystretch}{1.3}
\begin{tabular}{@{}lp{8cm}@{}}
\toprule
\textbf{Objekt-Typ} & \textbf{Beschreibung} \\
\midrule
Hauptdokument & Das eigentliche Dokument (PDF, Word, LaTeX) \\
Appendices & Ergänzende Materialien \\
Daten & Datasets, Code, Modelle \\
Präsentation & Slides, Poster \\
Executive Summary & Kurzfassung für Entscheider \\
Abstract & Kürzestfassung \\
Metadaten & Beschreibung, Keywords, Klassifikation \\
Begleitbriefe & Cover Letter, Submission Notes \\
\bottomrule
\end{tabular}
\end{table}


% =============================================================================
\section{Format-Ableitung: Von Dimensionen zum Dokumenttyp}
% =============================================================================

\subsection{Die Dokumenten-Format-Taxonomie}

\begin{table}[htbp]
\centering
\caption{EBF Dokumenten-Formate}
\renewcommand{\arraystretch}{1.4}
\small
\begin{tabular}{@{}lcccp{4cm}@{}}
\toprule
\textbf{Format} & \textbf{Seiten} & \textbf{Dauer} & \textbf{Qualität} & \textbf{Typischer Einsatz} \\
\midrule
\multicolumn{5}{@{}l}{\textit{\textbf{KURZ (1-5 Seiten)}}} \\
1-Pager & 1 & 2h & Q2-Q3 & Entscheidungsvorlage, Briefing \\
Executive Summary & 1-2 & 4h & Q3 & Zusammenfassung für Führung \\
Memo & 1-3 & 2-4h & Q2-Q3 & Interne Kommunikation \\
Abstract & 0.5-1 & 2h & Q3-Q4 & Konferenz, Journal \\
\addlinespace
\multicolumn{5}{@{}l}{\textit{\textbf{MITTEL (5-30 Seiten)}}} \\
Brief/Letter & 2-5 & 4h & Q3 & Formelle Korrespondenz \\
Proposal & 5-15 & 1-2T & Q3 & Projektantrag, Angebot \\
Report & 10-30 & 1-2W & Q3-Q4 & Analyse, Evaluation \\
Working Paper & 15-30 & 2-4W & Q2-Q3 & Forschung in Progress \\
Case Study & 10-20 & 1-2W & Q3 & Anwendungsbeispiel \\
\addlinespace
\multicolumn{5}{@{}l}{\textit{\textbf{LANG (30-100 Seiten)}}} \\
White Paper & 20-50 & 2-4W & Q3-Q4 & Position, Thought Leadership \\
Technical Report & 30-80 & 4-8W & Q3-Q4 & Detaillierte Analyse \\
Journal Article & 20-40 & 2-6M & Q4 & Peer-reviewed Publication \\
Thesis Chapter & 30-60 & 2-4M & Q4 & Dissertation \\
\addlinespace
\multicolumn{5}{@{}l}{\textit{\textbf{SEHR LANG (100+ Seiten)}}} \\
Monograph & 100-300 & 6-12M & Q4-Q5 & Buchkapitel, Einzelwerk \\
Handbook Chapter & 50-100 & 3-6M & Q4-Q5 & Referenzwerk \\
Textbook & 300-600 & 1-3J & Q5 & Lehrbuch \\
Comprehensive Treatise & 500+ & 2-5J & Q5 & Definitive Referenz \\
\bottomrule
\end{tabular}
\end{table}


\subsection{Format-Entscheidungsmatrix}

\begin{tcolorbox}[colback=green!5!white, colframe=green!75!black, 
    title=Format-Ableitung: Der Algorithmus]

\textbf{Schritt 1: Ziel + Leser → Basis-Kategorie}
\begin{center}
\begin{tabular}{@{}l|cccc@{}}
& Experten & Fachleute & Entscheider & Öffentlich \\
\hline
Informieren & Paper & Report & Summary & Artikel \\
Überzeugen & Paper & Proposal & Memo & Kampagne \\
Dokumentieren & Appendix & Manual & Protokoll & FAQ \\
Anleiten & Guide & Tutorial & Checklist & How-To \\
\end{tabular}
\end{center}

\textbf{Schritt 2: Scope + Evidenz → Länge}
\begin{itemize}[nosep]
\item Enger Scope + Anekdotisch → KURZ (1-5 S.)
\item Mittlerer Scope + Empirisch → MITTEL (5-30 S.)
\item Breiter Scope + Meta-analytisch → LANG (30-100 S.)
\item Umfassend + Formal → SEHR LANG (100+ S.)
\end{itemize}

\textbf{Schritt 3: Qualität + Vertraulichkeit → Finalformat}
\begin{itemize}[nosep]
\item Q4+ Öffentlich → Journal/Book
\item Q3+ Kunden → Report/White Paper
\item Q2+ Intern → Working Paper/Memo
\end{itemize}

\end{tcolorbox}


% =============================================================================
\section{EBF Spezifisch: Das Hauptdokument}
% =============================================================================

\begin{tcolorbox}[colback=blue!5!white, colframe=blue!75!black, 
    title=EBF Hauptdokument: Vollständige Spezifikation]

\textbf{1. ZIEL:} Theoretisieren + Integrieren
\begin{itemize}[nosep]
\item Unified Framework für Economic Rationality entwickeln
\item Bestehende Theorien als Spezialfälle subsumieren
\end{itemize}

\textbf{2. LESER:} Primär Experten, Sekundär Fachleute
\begin{itemize}[nosep]
\item Akademische Ökonomen (Behavioral, Institutional, Finance)
\item PhD-Studenten
\item Policy-Analysten mit Ökonomiehintergrund
\end{itemize}

\textbf{3. ANLASS:} Erkenntnisgewinn + Systematisierung
\begin{itemize}[nosep]
\item BEATRIX Research Group Forschungsprogramm
\item Vorbereitung für akademische Publikationen
\end{itemize}

\textbf{4. SCOPE (In-Scope):}
\begin{itemize}[nosep]
\item Komplementarität als Strukturprinzip
\item Kontext als endogene Variable
\item FEPSDE Welfare Framework (6 Dimensionen, 3 Kategorien)
\item Aggregationsebenen (L=0 bis L=$\infty$)
\item Awareness \& Willingness Funktionen
\item Falsifizierbare Vorhersagen
\end{itemize}

\textbf{5. OUT-OF-SCOPE:}
\begin{itemize}[nosep]
\item Vollständige mathematische Beweise (→ Appendices)
\item Implementierungsdetails (→ Separate Guides)
\item Softwarecode (→ GitHub Repository)
\item Einzelne Länder-/Branchen-Analysen (→ Case Studies)
\item Alternative Frameworks (nur erwähnt, nicht entwickelt)
\end{itemize}

\textbf{6. MINDESTANFORDERUNGEN:}
\begin{itemize}[nosep]
\item Klare Definitionen aller Kernkonzepte
\item Mathematische Kerngleichungen (Intuition + Form)
\item Verbindung zu bestehender Literatur (Nobel, Fehr, etc.)
\item Mindestens 6 falsifizierbare Vorhersagen
\item Worked Examples pro Aggregationsebene
\item Glossar aller Symbole
\end{itemize}

\textbf{7. QUALITÄTSSTANDARD:} Q4 (Publication Ready)
\begin{itemize}[nosep]
\item Peer-review Qualität für Hauptkapitel
\item Interne Review für alle Appendices
\end{itemize}

\textbf{8. EVIDENZ-ANFORDERUNG:}
\begin{itemize}[nosep]
\item Formal-mathematisch für Kernableitungen
\item Empirisch-kausal für Vorhersagen (wo möglich)
\item Theoretisch für Integration
\end{itemize}

\textbf{9. LEBENSZYKLUS:} Versioniert
\begin{itemize}[nosep]
\item Aktuell: Version 51 (Januar 2026)
\item Quarterly Minor Updates, Annual Major Updates
\end{itemize}

\textbf{10. VERTRAULICHKEIT:} Intern → Öffentlich (geplant)
\begin{itemize}[nosep]
\item Phase 1: FehrAdvice intern
\item Phase 2: Akademische Partner
\item Phase 3: Open Access
\end{itemize}

\textbf{11. ABHÄNGIGKEITEN:}
\begin{itemize}[nosep]
\item Completed: Core Theory, FEPSDE Matrix
\item In Progress: Awareness Function, Predictions
\item Pending: Empirical Validation
\end{itemize}

\textbf{12. LIEFEROBJEKTE:}
\begin{itemize}[nosep]
\item Hauptdokument (complementarity\_context\_v51.tex)
\item 49 Appendices
\item GitHub Repository mit Code
\item Präsentationen für verschiedene Audiences
\end{itemize}

\rule{\textwidth}{0.5pt}

\textbf{→ Abgeleitetes Format:} \textbf{Comprehensive Treatise / Textbook}
\begin{itemize}[nosep]
\item Umfang: 500+ Seiten (Haupt) + 750+ Seiten (Appendices)
\item Struktur: 7 Teile, 19 Kapitel, ~50 Subkapitel
\item Qualität: Q4-Q5
\item Timeline: Multi-Jahr Projekt
\end{itemize}

\end{tcolorbox}


% =============================================================================
\section{Quick Reference: Spezifikations-Template}
% =============================================================================

\begin{tcolorbox}[colback=gray!10!white, colframe=black, 
    title=DOCUMENT SPECIFICATION CARD (Template)]

\textbf{Dokument-Titel:} \_\_\_\_\_\_\_\_\_\_\_\_\_\_\_\_\_\_\_\_\_\_\_\_\_\_\_\_

\textbf{Version/Datum:} \_\_\_\_\_\_\_\_\_\_\_\_

\rule{\textwidth}{0.3pt}

\textbf{1. ZIEL:} $\square$ Informieren $\square$ Überzeugen $\square$ Dokumentieren 
$\square$ Anleiten $\square$ Entscheiden $\square$ Theoretisieren

\textbf{2. LESER:} $\square$ Experten $\square$ Fachleute $\square$ Generalisten 
$\square$ Entscheider $\square$ Studenten $\square$ Öffentlichkeit

\textbf{3. ANLASS:} $\square$ Anfrage $\square$ Deadline $\square$ Meilenstein 
$\square$ Erkenntnisgewinn $\square$ Systematisierung $\square$ Pflicht

\rule{\textwidth}{0.3pt}

\textbf{4. SCOPE (Was wird behandelt):}

\_\_\_\_\_\_\_\_\_\_\_\_\_\_\_\_\_\_\_\_\_\_\_\_\_\_\_\_\_\_\_\_\_\_\_\_\_\_\_\_

\_\_\_\_\_\_\_\_\_\_\_\_\_\_\_\_\_\_\_\_\_\_\_\_\_\_\_\_\_\_\_\_\_\_\_\_\_\_\_\_

\textbf{5. OUT-OF-SCOPE (Was wird NICHT behandelt):}

\_\_\_\_\_\_\_\_\_\_\_\_\_\_\_\_\_\_\_\_\_\_\_\_\_\_\_\_\_\_\_\_\_\_\_\_\_\_\_\_

\_\_\_\_\_\_\_\_\_\_\_\_\_\_\_\_\_\_\_\_\_\_\_\_\_\_\_\_\_\_\_\_\_\_\_\_\_\_\_\_

\textbf{6. MINDESTANFORDERUNGEN (Must-Haves):}

\_\_\_\_\_\_\_\_\_\_\_\_\_\_\_\_\_\_\_\_\_\_\_\_\_\_\_\_\_\_\_\_\_\_\_\_\_\_\_\_

\_\_\_\_\_\_\_\_\_\_\_\_\_\_\_\_\_\_\_\_\_\_\_\_\_\_\_\_\_\_\_\_\_\_\_\_\_\_\_\_

\rule{\textwidth}{0.3pt}

\textbf{7. QUALITÄT:} $\square$ Q1 Draft $\square$ Q2 Working $\square$ Q3 Professional 
$\square$ Q4 Publication $\square$ Q5 Reference

\textbf{8. EVIDENZ:} $\square$ Anekdotisch $\square$ Expert Opinion $\square$ Theoretisch 
$\square$ Empirisch $\square$ Meta-analytisch $\square$ Formal

\textbf{9. LEBENSZYKLUS:} $\square$ Einmalig $\square$ Versioniert $\square$ Living Document

\textbf{10. VERTRAULICHKEIT:} $\square$ Öffentlich $\square$ Kunden $\square$ Intern 
$\square$ Restricted $\square$ Confidential

\rule{\textwidth}{0.3pt}

\textbf{11. ABHÄNGIGKEITEN:}

\_\_\_\_\_\_\_\_\_\_\_\_\_\_\_\_\_\_\_\_\_\_\_\_\_\_\_\_\_\_\_\_\_\_\_\_\_\_\_\_

\textbf{12. LIEFEROBJEKTE:}

\_\_\_\_\_\_\_\_\_\_\_\_\_\_\_\_\_\_\_\_\_\_\_\_\_\_\_\_\_\_\_\_\_\_\_\_\_\_\_\_

\rule{\textwidth}{0.3pt}

\textbf{→ ABGELEITETES FORMAT:} \_\_\_\_\_\_\_\_\_\_\_\_\_\_\_\_\_\_\_\_\_\_

\textbf{→ GESCHÄTZTE LÄNGE:} \_\_\_\_\_ Seiten

\textbf{→ GESCHÄTZTER AUFWAND:} \_\_\_\_\_

\end{tcolorbox}


% =============================================================================
\section{Anhang: Format-Entscheidungsbaum (Flowchart)}
% =============================================================================

\begin{verbatim}
START
  │
  ▼
┌─────────────────────────────────────┐
│ Wie viele Leser?                    │
├─────────────────────────────────────┤
│ 1-5 Personen → Memo/Email           │
│ 6-50 Personen → Report              │
│ 50+ Personen → Publication          │
└─────────────────────────────────────┘
  │
  ▼
┌─────────────────────────────────────┐
│ Wie viel Zeit hat der Leser?        │
├─────────────────────────────────────┤
│ <5 Minuten → 1-Pager                │
│ 5-30 Minuten → Summary/Memo         │
│ 30-120 Minuten → Report             │
│ 2+ Stunden → Paper/Book             │
└─────────────────────────────────────┘
  │
  ▼
┌─────────────────────────────────────┐
│ Wie komplex ist das Thema?          │
├─────────────────────────────────────┤
│ Einfach (1 Hauptaussage) → 1-Pager  │
│ Mittel (3-5 Punkte) → Report        │
│ Komplex (System) → Paper            │
│ Sehr komplex (Theorie) → Book       │
└─────────────────────────────────────┘
  │
  ▼
┌─────────────────────────────────────┐
│ Welche Evidenz wird erwartet?       │
├─────────────────────────────────────┤
│ Keine/Meinung → Memo                │
│ Belege/Daten → Report               │
│ Peer-review → Journal Paper         │
│ Exhaustiv → Monograph               │
└─────────────────────────────────────┘
  │
  ▼
FORMAT BESTIMMT
\end{verbatim}


\end{document}
